\chapter*{Resumo} \label{chap:Resumo}

Muitos dos sinais tratados por sensores inteligentes são não periódicos. Mantêm-se quase constantes durante longos
intervalos e só ocasionalmente mudam, como acontece com uma leitura de
temperatura ao longo do tempo. Um conversor analógico-digital síncrono
convencional amostra uma entrada deste tipo a uma taxa de relógio fixa, pelo que,
o seu banco de comparadores e a rede de distribuição de relógio continuam a
comutar e a dissipar potência.

Esta dissertação, projeta um conversor \textit{Flash} assíncrono que remove o relógio
global e converte por cruzamento de níveis (\textit{level-crossing}), conduzido pelo
próprio sinal. Desta forma a potência dinâmica acompanha a atividade da
entrada e reduz-se para um nível estático quando a entrada está em repouso. Remover o
relógio elimina também a rede de distribuição de alta velocidade, que representa
uma grande parte do orçamento de potência de um \textit{Flash} síncrono.

O conversor é projetado num processo CMOS de 16~nm a 0.8~V, com resolução de
4 bits e um banco de quinze múltiplos de comparação. Um amplificador operacional de
transcondutância complementar de cinco transístores, usado como seguidor de
tensão de ganho unitário, reproduz a entrada em quase toda a gama de alimentação.
 Cada múltiplo é um comparador de corrente cujo nível de decisão é definido pelo tamanho dos seus
transístores, pelo que os quinze limiares resultam do dimensionamento e não de
um conjunto de referências. Como esses limiares dependentes do dimensionamento
variam com o processo de fabrico, uma estratégia de calibração por
polarização do substrato (\textit{body bias}), implementada como modelo comportamental
em Verilog-A, ajusta as tensões de bulk de cada comparador para repor os pontos de
disparo ao longo das variações de processo. Realizado na Synopsys Portugal, o
trabalho engloba a análise, o projeto ao nível de esquemático e a validação por
co-simulação analógica.

O conversor é comparado com uma referência Flash síncrona da mesma resolução e
tecnologia, caracterizada nas mesmas condições. Com a mesma carga de conversões
por segundo, o núcleo assíncrono consome cerca de um nono da potência da
referência, sendo que, a maior parte dessa diferença a distribuição de relógio que o
projeto assíncrono não tem, e a poupança aumenta para entradas não periódicas à
medida que o conversor se aproxima do seu nível de repouso.